\section{Konu 4: Çok Değişkenli Fonksiyonlar, Limit ve Süreklilik}

Bu bölümde derste (15.10 ve 22.10 defter notları, 23.10 ve 08.11 tahta fotoğrafları) çözülen bütün örnekler, işlenme
sırasıyla ve adım adım çözümleriyle verilmiştir. Aynı örnek birden çok fotoğrafta geçiyorsa bir kez yazılmış, bütün
kaynak fotoğraflar belirtilmiştir. Sonda kaynak kitaplardaki tipik alıştırmalara benzer ek alıştırmalar vardır.

\begin{bilgi}[Çalışma soruları hakkında]
Bu konu için tahtada ayrıca ``çalışma sorusu / problemler'' listesi verilmemiştir. \foto{IMG-20261015-WA0040} ve
\foto{IMG-20261023-WA0042}, \foto{IMG-20261023-WA0043} kenarlarında görünen problem listesi (Frenet üçlüsü, vektör
integralleri) Konu 3'e aittir ve orada çözülmüştür.
\end{bilgi}

% =====================================================================
\subsection{Fonksiyon kavramı, tanım ve görüntü kümeleri}

\begin{soru}[Soru 1 \foto{IMG-20261015-WA0040}, \foto{IMG-20261023-WA0043}]
Aşağıdaki büyüklükleri fonksiyon olarak yazınız ve kaç değişkenli olduklarını belirtiniz:
(a) tabanı $a$, yüksekliği $h$ olan üçgenin alanı; (b) $A$ anaparanın $t$ faiz oranıyla $n$ yıllık basit faizi;
(c) kenarları $x,y$ olan dikdörtgenin alanı; (d) boyutları $x,y,h$ olan kutunun hacmi.
\end{soru}
\begin{cevap}
\begin{description}
  \item[(a)] $A(a,h)=\dfrac{a\cdot h}{2}$: iki bağımsız değişken ($a,h$), \textbf{iki değişkenli}. Fiziksel anlamı
  gereği tanım kümesi $a>0,\ h>0$ (birinci bölgenin içi).
  \item[(b)] $f(A,n,t)=\dfrac{A\cdot n\cdot t}{100}$ ($t$ yüzde olarak): \textbf{üç değişkenli}.
  \item[(c)] $f(x,y)=x\cdot y$: \textbf{iki değişkenli}.
  \item[(d)] $f(x,y,h)=x\cdot y\cdot h$: \textbf{üç değişkenli}.
\end{description}
Genel olarak $D\subset\mathbb R^2$ ise $z=f(x,y)$, $D\subset\mathbb R^3$ ise $w=f(x,y,z)$ yazılır; her girdi noktasına
\emph{tek} bir reel sayı karşılık gelir.
\end{cevap}

\begin{soru}[Soru 2 \foto{IMG-20261015-WA0041}, \foto{IMG-20261023-WA0042}]
$f(x,y)=\sqrt{4-x^2-y^2}$ eşitliği ile tanımlanan $f$ reel değerli fonksiyonunun tanım kümesini analitik düzlemde
gösteriniz.
\end{soru}
\begin{cevap}[Çözüm (adım adım)]
\textbf{1. Koşulu yaz.} Karekökün içi negatif olamaz: $4-x^2-y^2\ge0$.

\textbf{2. Düzenle.} $x^2+y^2\le4=2^2$.

\textbf{3. Geometrik yorum.} $x^2+y^2=4$ merkezi orijin, yarıçapı $2$ olan çemberdir; ``$\le$'' çemberin içi ve kendisi
demektir:
\[
T_f=\{(x,y): x^2+y^2\le4\}\quad\text{(kapalı disk, sınır çember dahil)}.
\]
\textbf{4. Çizim.} Çember sürekli çizgiyle çizilir, içi taranır.

\textbf{Ek (görüntü kümesi).} Disk içinde $0\le4-x^2-y^2\le4$ olduğundan $0\le f\le2$; orijinde $2$, çember üzerinde
$0$ değeri alınır ve arada her değer alınır: $G_f=[0,2]$. Grafik, yarıçapı $2$ olan kürenin üst yarısıdır.
\end{cevap}

\begin{soru}[Soru 3 \foto{IMG-20261015-WA0041}, \foto{IMG-20261023-WA0042}]
$f(x,y)=\arccos(x-y)$ fonksiyonunun tanım ve görüntü kümesi nedir?
\end{soru}
\begin{cevap}[Çözüm (adım adım)]
\textbf{1. Hatırlatma.} $\cos:\mathbb R\to[-1,1]$; $\arccos:[-1,1]\to[0,\pi]$. Yani $\arccos$'un girdisi $[-1,1]$
aralığında olmalıdır.

\textbf{2. Koşul.} $-1\le x-y\le1$.

\textbf{3. İki eşitsizliğe ayır.}
\[
x-y\ge-1\ \Rightarrow\ y\le x+1, \qquad x-y\le1\ \Rightarrow\ y\ge x-1 .
\]
\textbf{4. Tanım kümesi.}
\[
T_f=\{(x,y): x-1\le y\le x+1\}:
\]
eğimi $1$ olan paralel $y=x-1$ ve $y=x+1$ doğruları arasındaki şerit; sınır doğruları dahil.

\textbf{5. Görüntü kümesi.} Şerit üzerinde $u=x-y$ ifadesi $[-1,1]$'in her değerini alır (ör. $(x,y)=(u,0)$).
$\arccos u$ bu aralıkta $\arccos(-1)=\pi$'den $\arccos1=0$'a sürekli ve azalan biçimde iner. Dolayısıyla
\[
G_f=[0,\pi].
\]
\emph{Not:} Tahtadaki $\arccos u$ çizimi kabataslaktır ve etiketi \okunamadi{``arccos4'' gibi görünüyor; büyük olasılıkla
``$\arccos u$''}; gerçek eğri $(0,\pi/2)$'den geçer.
\end{cevap}

\begin{soru}[Soru 4 \foto{IMG-20261015-WA0042}, \foto{IMG-20261023-WA0044}]
$f(x,y)=\sqrt{xy}+\log(1-x^2-y^2)$ şeklinde tanımlı $f$ fonksiyonunun tanım kümesi nedir?
\end{soru}
\begin{cevap}[Çözüm (adım adım)]
\textbf{1. Kök koşulu.} $xy\ge0$: çarpım negatif değilse $x$ ve $y$ aynı işaretli ya da biri sıfırdır:
$(x\ge0$ ve $y\ge0)$ veya $(x\le0$ ve $y\le0)$. Bu, I. ve III. bölgeler (eksenler dahil).

\textbf{2. Logaritma koşulu.} $1-x^2-y^2>0\Rightarrow x^2+y^2<1$: birim çemberin içi, çember \textbf{hariç} (kesin
eşitsizlik).

\textbf{3. Kesişim.}
\[
T_f=\{(x,y): xy\ge0,\ x^2+y^2<1\}.
\]
Birim dairenin I. ve III. bölgedeki çeyrekleri; eksen parçaları dahil, çember hariç (çizimde çember kesikli).
Orijin de dahildir: $\sqrt0+\log1=0$. Logaritmanın tabanı tanım kümesini etkilemez.
\end{cevap}

\begin{soru}[Soru 5 \foto{IMG-20261022-WA0053}, \foto{IMG-20261108-WA0088}]
Aşağıdaki fonksiyonların tanım kümelerini bulunuz.
(a) $f(x,y)=\ln(x^2+y^2+4)$ \quad (b) $f(x,y)=\sqrt{xy-x^4}+xy$ \quad (c) $f(x,y)=\sqrt{x^2-y}+x^2+y^2+1$
\end{soru}
\begin{cevap}[Çözüm (adım adım)]
\begin{description}
  \item[(a)] Logaritma koşulu: $x^2+y^2+4>0$. $x^2+y^2\ge0$ olduğundan sol taraf her zaman en az $4$'tür; koşul her
  $(x,y)$ için doğrudur. $T_f=\mathbb R^2$ (bütün düzlem).
  \item[(b)] Kök koşulu: $xy-x^4\ge0$. Ortak çarpan $x$: $x(y-x^3)\ge0$. Çarpımın işaret tablosu:
  \begin{itemize}
    \item (1) $x\ge0$ ve $y-x^3\ge0$, yani $x\ge0$ ve $y\ge x^3$: sağ yarı düzlemde $y=x^3$ eğrisinin üstü;
    \item (2) $x\le0$ ve $y-x^3\le0$, yani $x\le0$ ve $y\le x^3$: sol yarı düzlemde $y=x^3$ eğrisinin altı.
  \end{itemize}
  $T_f$ = (1) $\cup$ (2); sınırlar ($y=x^3$ eğrisi ve $y$ ekseninin tamamı) dahildir. Kontrol: $(1,2)$ için
  $2-1=1\ge0$ \checkmark; $(1,0)$ için $0-1<0$, dışarıda \checkmark; $(-1,-2)$ için $2-1=1\ge0$ \checkmark.
  \item[(c)] Kök koşulu: $x^2-y\ge0\Rightarrow y\le x^2$. Polinom kısmı her yerde tanımlı. $T_f$: $y=x^2$ parabolünün
  altı ve kendisi. Kontrol: $(0,-1)$: $0+1=1\ge0$ \checkmark; $(0,1)$: $-1<0$, dışarıda \checkmark.
\end{description}
\end{cevap}

% =====================================================================
\subsection{Seviye eğrileri ve grafik}

\begin{soru}[Soru 6 \foto{IMG-20261015-WA0042}, \foto{IMG-20261015-WA0043}, \foto{IMG-20261023-WA0044}]
$f(x,y)=x^2+y^2$ fonksiyonunun bazı seviye eğrilerini bulunuz. Bundan yararlanarak fonksiyon grafiğini çiziniz.
\end{soru}
\begin{cevap}[Çözüm (adım adım)]
\textbf{1. Seviye eğrisi denklemi.} $f(x,y)=c\ \Rightarrow\ x^2+y^2=c$.

\textbf{2. $c$'ye göre durum.}
\begin{itemize}
  \item $c<0$: $x^2+y^2\ge0$ olduğundan çözüm yok (boş küme).
  \item $c=0$: yalnızca $(0,0)$.
  \item $c>0$: merkezi orijin, yarıçapı $\sqrt c$ olan çember. Örneğin $c=1$: birim çember; $c=4$: yarıçapı $2$;
        $c=9$: yarıçapı $3$.
\end{itemize}
\textbf{3. Grafiği kurma.} Yüzeyin $z=c$ yüksekliğindeki kesiti yarıçapı $\sqrt c$ olan yatay bir çemberdir. Yükseklik
arttıkça çemberler büyür; $x=0$ düzlemiyle kesit $z=y^2$, $y=0$ düzlemiyle kesit $z=x^2$ parabolüdür. Bu bilgiler
birleşince orijinden yukarı açılan kâse biçimli \textbf{paraboloid} $z=x^2+y^2$ elde edilir.

\textbf{Gösterim notu.} Tahtada ``$x^2+y^2=c^2$ alındığında yarıçapı $c$ olan çember'' yazılmıştır; bu, seviye değeri
$c^2$ seçildiğinde doğrudur. $f=c$ yazımıyla doğrusu yarıçapın $\sqrt c$ olmasıdır; tahtadaki şekildeki $c=1,4,9$
etiketleri (yarıçap $1,2,3$) de bunu gösterir. Ayrıca tahtadaki paraboloid düz kenarlı (koni gibi) çizilmiştir;
gerçek yüzeyin kenarları eğridir.
\end{cevap}

% =====================================================================
\subsection{Limit hesabı ve limitin yokluğu}

\begin{soru}[Soru 7 \foto{IMG-20261022-WA0058}, \foto{IMG-20261108-WA0094}]
$\displaystyle\lim_{(x,y)\to(0,0)}\frac{2xy}{x^2+y^2}$ limitini inceleyiniz.
\end{soru}
\begin{cevap}[Çözüm (adım adım)]
\textbf{1. Doğrudan yerine koyma} $\frac00$ belirsizliği verir; fonksiyon orijinde tanımsız.

\textbf{2. Kutupsal koordinatlar.} $x=r\cos\theta$, $y=r\sin\theta$; $(x,y)\to(0,0)\Leftrightarrow r\to0$:
\[
\frac{2xy}{x^2+y^2}=\frac{2r^2\cos\theta\sin\theta}{r^2(\cos^2\theta+\sin^2\theta)}=2\cos\theta\sin\theta=\sin2\theta .
\]
\textbf{3. Yorum.} Sonuç $r$'den bağımsız ama $\theta$'ya bağlı: $\theta=0$ doğrultusunda $0$, $\theta=\pi/4$
doğrultusunda $1$, $\theta=3\pi/4$ doğrultusunda $-1$. Farklı doğrultularda farklı değer $\Rightarrow$ \textbf{limit
yoktur}.

\textbf{4. Doğrularla doğrulama.} $y=mx$:
\[
\frac{2x\cdot mx}{x^2+m^2x^2}=\frac{2m}{1+m^2};\qquad m=0\Rightarrow0,\quad m=1\Rightarrow1 .
\]
İki farklı yol, iki farklı değer: limit yoktur.
\end{cevap}

\begin{soru}[Soru 8 \foto{IMG-20261022-WA0053}]
$P(x,y)=5x^2y^2+4xy^3-8xy^2+3$ için $\displaystyle\lim_{(x,y)\to(1,2)}P(x,y)$ değerini bulunuz.
\end{soru}
\begin{cevap}
Polinomlar için $\lim_{(x,y)\to(a,b)}x^my^n=a^mb^n$ ve limit kurallarından $\lim P(x,y)=P(a,b)$ (doğrudan yerine koyma):
\[
P(1,2)=5\cdot1\cdot4+4\cdot1\cdot8-8\cdot1\cdot4+3=20+32-32+3=\boxed{23}.
\]
\end{cevap}

\begin{soru}[Soru 9 \foto{IMG-20261022-WA0053}, \foto{IMG-20261022-WA0056}, \foto{IMG-20261108-WA0088}, \foto{IMG-20261108-WA0089}]
Aşağıdaki limitleri hesaplayınız.
\begin{description}
  \item[(a)] $\displaystyle\lim_{(x,y)\to(0,0)}\frac{2x^2+y^2+1}{x^2+y^2+3}$ \qquad
  (b) $\displaystyle\lim_{(x,y)\to(0,0)}\frac{2-\sqrt{xy+4}}{xy}$
  \item[(c)] $\displaystyle\lim_{(x,y)\to(0,0)}\frac{\sin(xy)}{x}$ \qquad
  (d) $\displaystyle\lim_{(x,y)\to(0,\ln2)}e^{x-y}$
  \item[(e)] $\displaystyle\lim_{(x,y)\to(1,0)}\frac{x\sin y}{x^2+1}$ \qquad
  (f) $\displaystyle\lim_{(x,y)\to(0,0)}\frac{e^y\sin x}{x}$
\end{description}
\end{soru}
\begin{cevap}[Çözüm (adım adım)]
\begin{description}
  \item[(a)] Payda $(0,0)$'da $3\neq0$. Rasyonel fonksiyon, doğrudan yerine koyma:
  $\dfrac{0+0+1}{0+0+3}=\dfrac13$.
  \item[(b)] \textbf{1.} Yerine koyma: $\frac{2-\sqrt4}{0}=\frac00$, belirsiz.
  \textbf{2.} Payın eşleniği $2+\sqrt{xy+4}$ ile pay ve paydayı çarp ($xy\neq0$ noktalarında çalışıyoruz):
  \[
  \frac{(2-\sqrt{xy+4})(2+\sqrt{xy+4})}{xy\,(2+\sqrt{xy+4})}=\frac{4-(xy+4)}{xy\,(2+\sqrt{xy+4})}
  =\frac{-xy}{xy\,(2+\sqrt{xy+4})}=\frac{-1}{2+\sqrt{xy+4}} .
  \]
  \textbf{3.} Artık payda sıfırdan farklı; yerine koy: $\dfrac{-1}{2+\sqrt4}=-\dfrac14$.
  \item[(c)] \textbf{1.} $\frac00$ belirsizliği. \textbf{2.} Pay ve paydayı $y$ ile çarpıp $u=xy\to0$ kullan:
  \[
  \frac{\sin(xy)}{x}=y\cdot\frac{\sin(xy)}{xy}\ \Rightarrow\
  \lim\frac{\sin(xy)}{xy}\cdot\lim y=1\cdot0=0 .
  \]
  \textbf{Titiz alternatif:} $|\sin u|\le|u|$ ile $\left|\frac{\sin(xy)}{x}\right|\le\frac{|x||y|}{|x|}=|y|\to0$;
  sıkıştırma ile limit $0$ (bu yol $y=0$ noktalarında tanımsız bölme yaratmaz).
  \item[(d)] $e^{x-y}=e^x\cdot e^{-y}$:
  \[
  \lim_{x\to0}e^x\cdot\lim_{y\to\ln2}e^{-y}=1\cdot e^{-\ln2}=1\cdot\frac12=\frac12 .
  \]
  \emph{Düzeltme:} 22.10 defter notunda (\foto{IMG-20261022-WA0056}) ikinci çarpan $e^{y}$ yazılmış; doğrusu
  $e^{-y}$'dir (08.11 tahtasında, \foto{IMG-20261108-WA0089}, doğru yazılmış). $e^y$ ile hesaplansaydı $2$ çıkardı.
  Sonuç $\tfrac12$ doğrudur.
  \item[(e)] Payda $x^2+1\ge1$ asla sıfır olmaz; değişkenlere ayır:
  \[
  \lim_{x\to1}\frac{x}{x^2+1}\cdot\lim_{y\to0}\sin y=\frac{1}{2}\cdot0=0 .
  \]
  \item[(f)] $\dfrac{e^y\sin x}{x}=e^y\cdot\dfrac{\sin x}{x}$; iki tek değişkenli limit de var:
  \[
  \lim_{x\to0}\frac{\sin x}{x}\cdot\lim_{y\to0}e^y=1\cdot1=1 .
  \]
  (22.10 defterinde yarım kalan bu şıkkın çözümü 08.11 tahtasından alınmıştır.)
\end{description}
\end{cevap}

\begin{soru}[Soru 10 \foto{IMG-20261022-WA0056}, \foto{IMG-20261022-WA0054}, \foto{IMG-20261108-WA0089}]
Aşağıdaki fonksiyonların $(0,0)$ noktasında limite sahip olmadığını gösteriniz.
\begin{description}
  \item[(a)] $f(x,y)=\dfrac{x^4-y^2}{x^4+y^2}$ \qquad (b) $f(x,y)=\dfrac{x-y}{x+y}$
  \item[(c)] $f(x,y)=\dfrac{x}{\sqrt{x^2+y^2}}$ \qquad (d) $f(x,y)=\dfrac{x^2+y^2}{y^2}$
\end{description}
\end{soru}
\begin{cevap}[Çözüm (adım adım)]
Strateji: iki farklı yol bulup farklı limit değerleri elde etmek.
\begin{description}
  \item[(a)] Paydada $x^4$ ve $y^2$ olduğu için $y=ax^2$ parabolleri $x^4$ ile $y^2$'yi aynı mertebeye getirir:
  \[
  \lim_{x\to0}\frac{x^4-a^2x^4}{x^4+a^2x^4}=\frac{1-a^2}{1+a^2}.
  \]
  $a=0$ ($x$ ekseni) $\Rightarrow1$; $a=1$ ($y=x^2$) $\Rightarrow0$. Farklı $\Rightarrow$ limit yok.
  \item[(b)] $y=mx$ ($m\neq-1$, çünkü $y=-x$ üzerinde payda sıfır):
  \[
  \lim_{x\to0}\frac{x-mx}{x+mx}=\lim_{x\to0}\frac{x(1-m)}{x(1+m)}=\frac{1-m}{1+m}.
  \]
  $m=0\Rightarrow1$, $m=1\Rightarrow0$. Limit yok.
  \item[(c)] Kutupsal ($r>0$ olduğundan $\sqrt{r^2}=r$):
  \[
  \frac{r\cos\theta}{\sqrt{r^2(\cos^2\theta+\sin^2\theta)}}=\frac{r\cos\theta}{r}=\cos\theta .
  \]
  $\theta=0$'da $1$, $\theta=\pi/2$'de $0$, $\theta=\pi$'de $-1$; $\theta$'ya bağlı $\Rightarrow$ limit yok.
  \item[(d)] $y=mx$ ($m\neq0$; $y=0$ üzerinde $f$ tanımsız):
  \[
  \lim_{x\to0}\frac{x^2+m^2x^2}{m^2x^2}=\frac{x^2(1+m^2)}{x^2m^2}=\frac{1+m^2}{m^2}.
  \]
  $m=1\Rightarrow2$, $m=2\Rightarrow\frac54$. Limit yok.
\end{description}
\end{cevap}

% =====================================================================
\subsection{Süreklilik}

\begin{soru}[Soru 11 \foto{IMG-20261022-WA0054}, \foto{IMG-20261108-WA0090}]
\[
f(x,y)=\begin{cases}\dfrac{\sin(xy^2)}{x^2+y^2}, & (x,y)\neq(0,0)\\[6pt] 0, & (x,y)=(0,0)\end{cases}
\]
fonksiyonunun her yerde sürekli olduğunu gösteriniz.
\end{soru}
\begin{cevap}[Çözüm (adım adım)]
Süreklilik tanımı: $\lim_{(x,y)\to(a,b)}f(x,y)=f(a,b)$ (\foto{IMG-20261022-WA0053}, \foto{IMG-20261108-WA0088}).

\textbf{1. $(x,y)\neq(0,0)$.} $\sin(xy^2)$, sürekli $\sin$ ile $xy^2$ polinomunun bileşkesi olduğundan süreklidir;
$x^2+y^2$ polinomdur ve bu noktalarda payda sıfır değildir. Bölüm sürekli $\Rightarrow f$ süreklidir.

\textbf{2. $(0,0)$ noktası.} $f(0,0)=0$ tanımlı. Limit:
\[
\lim_{(x,y)\to(0,0)}\frac{\sin(xy^2)}{x^2+y^2}=\lim\frac{\sin(xy^2)}{xy^2}\cdot\lim\frac{xy^2}{x^2+y^2}.
\]
Birinci çarpan: $u=xy^2\to0$ ve $\frac{\sin u}{u}\to1$. İkinci çarpan:
\begin{itemize}
  \item[(1)] $y=mx$: $\dfrac{x\cdot m^2x^2}{x^2+m^2x^2}=\dfrac{m^2x}{1+m^2}\to0$.
  \item[(2)] Kutupsal: $\dfrac{r\cos\theta\cdot r^2\sin^2\theta}{r^2}=r\cos\theta\sin^2\theta$ ve
  $\boxed{|r\cos\theta\sin^2\theta|\le r}\to0$.
\end{itemize}
Limit $=1\cdot0=0=f(0,0)$; $f$ orijinde de süreklidir. \textbf{Sonuç:} $f$ her yerde süreklidir.

\textbf{Titizlik notu.} (1)'deki doğru yolu tek başına kanıt değildir (doğrularda $0$ çıkıp eğride farklı değer çıkan
fonksiyonlar vardır); kanıtı $\theta$'dan bağımsız (2) sınırı tamamlar. Ayrıca $\frac{\sin(xy^2)}{xy^2}$ eksenler
üzerinde tanımsızdır. Bölme yapmadan tek satırlık kanıt: $|\sin u|\le|u|$ ile
\[
|f(x,y)|\le\frac{|x|\,y^2}{x^2+y^2}\le|x|\to0 .
\]
\end{cevap}

% =====================================================================
\subsection{Ek alıştırmalar}

\begin{soru}[\ekal Ek alıştırma 1]
$f(x,y)=\sqrt{1-x^2}+\sqrt{1-y^2}$ fonksiyonunun tanım ve görüntü kümesini bulunuz.
\end{soru}
\begin{cevap}
Kök koşulları: $1-x^2\ge0\Rightarrow-1\le x\le1$ ve $1-y^2\ge0\Rightarrow-1\le y\le1$.
$T_f=[-1,1]\times[-1,1]$: köşeleri $(\pm1,\pm1)$ olan kapalı kare.

Her terim $[0,1]$ aralığında değer alır; toplam $[0,2]$ içindedir. $f(\pm1,\pm1)=0$, $f(0,0)=2$; $y=0$ alıp
$x$'i $1$'den $0$'a götürdüğümüzde $f=\sqrt{1-x^2}+1$ değeri $1$'den $2$'ye, $x=1$ alıp $y$'yi değiştirdiğimizde
$0$'dan $1$'e sürekli değişir. Böylece $G_f=[0,2]$.
\end{cevap}

\begin{soru}[\ekal Ek alıştırma 2]
$f(x,y)=x^2-y^2$ fonksiyonunun seviye eğrilerini $c=-1,0,1$ için bulunuz ve grafiğin biçimini yorumlayınız.
\end{soru}
\begin{cevap}
\begin{itemize}
  \item $c=1$: $x^2-y^2=1$; tepe noktaları $(\pm1,0)$ olan, $x$ ekseni boyunca açılan hiperbol.
  \item $c=0$: $x^2=y^2\Rightarrow y=\pm x$; orijinde kesişen iki doğru.
  \item $c=-1$: $y^2-x^2=1$; tepe noktaları $(0,\pm1)$ olan, $y$ ekseni boyunca açılan hiperbol.
\end{itemize}
Grafik $z=x^2-y^2$, $x$ yönünde yukarı ($z=x^2$), $y$ yönünde aşağı ($z=-y^2$) açılan bir \textbf{eyer yüzeyi}
(hiperbolik paraboloid) belirtir.
\end{cevap}

\begin{soru}[\ekal Ek alıştırma 3]
$\varepsilon$--$\delta$ tanımıyla $\displaystyle\lim_{(x,y)\to(1,2)}(2x+3y)=8$ olduğunu gösteriniz.
\end{soru}
\begin{cevap}
$\varepsilon>0$ verilsin ve $d=\sqrt{(x-1)^2+(y-2)^2}$ olsun. $|x-1|\le d$ ve $|y-2|\le d$ olduğundan
\[
|2x+3y-8|=|2(x-1)+3(y-2)|\le2|x-1|+3|y-2|\le5d .
\]
$\delta=\varepsilon/5$ seçilirse $0<d<\delta$ iken $|2x+3y-8|<5\delta=\varepsilon$ olur. Limit $8$'dir.
\end{cevap}

\begin{soru}[\ekal Ek alıştırma 4]
Limitleri hesaplayınız:
(a) $\displaystyle\lim_{(x,y)\to(1,1)}\frac{x^3-y^3}{x-y}$ \qquad
(b) $\displaystyle\lim_{(x,y)\to(0,0)}\frac{x^2+y^2}{\sqrt{x^2+y^2+1}-1}$
\end{soru}
\begin{cevap}
\begin{description}
  \item[(a)] $\frac00$ belirsizliği ($x\neq y$). $x^3-y^3=(x-y)(x^2+xy+y^2)$:
  \[
  \lim_{(x,y)\to(1,1)}(x^2+xy+y^2)=1+1+1=3 .
  \]
  \item[(b)] $\frac00$ belirsizliği. Paydanın eşleniği $\sqrt{x^2+y^2+1}+1$ ile çarp:
  \[
  \frac{(x^2+y^2)\big(\sqrt{x^2+y^2+1}+1\big)}{(x^2+y^2+1)-1}=\sqrt{x^2+y^2+1}+1\ \longrightarrow\ 1+1=2 .
  \]
\end{description}
\end{cevap}

\begin{soru}[\ekal Ek alıştırma 5]
Aşağıdaki limitlerin olmadığını gösteriniz:
(a) $\displaystyle\lim_{(x,y)\to(0,0)}\frac{x^2-y^2}{x^2+y^2}$ \qquad
(b) $\displaystyle\lim_{(x,y)\to(0,0)}\frac{xy^2}{x^2+y^4}$
\end{soru}
\begin{cevap}
\begin{description}
  \item[(a)] $x$ ekseni ($y=0$): $\frac{x^2}{x^2}=1$. $y$ ekseni ($x=0$): $\frac{-y^2}{y^2}=-1$. Farklı; limit yok.
  \item[(b)] Doğrular $y=mx$: $\dfrac{m^2x^3}{x^2+m^4x^4}=\dfrac{m^2x}{1+m^4x^2}\to0$; eksenlerde de $0$. Ama $x=y^2$
  parabolü boyunca: $\dfrac{y^2\cdot y^2}{y^4+y^4}=\dfrac12$. Doğrularda $0$, parabolde $\frac12$; limit yok.
  (Bütün doğrularda aynı değerin çıkması limitin varlığı için yetmez!)
\end{description}
\end{cevap}

\begin{soru}[\ekal Ek alıştırma 6]
$\displaystyle\lim_{(x,y)\to(0,0)}\frac{x^4+y^4}{x^2+y^2}$ limitinin var olduğunu gösterip değerini bulunuz.
\end{soru}
\begin{cevap}
$x^4+y^4\le x^4+2x^2y^2+y^4=(x^2+y^2)^2$ olduğundan $(x,y)\neq(0,0)$ için
\[
0\le\frac{x^4+y^4}{x^2+y^2}\le\frac{(x^2+y^2)^2}{x^2+y^2}=x^2+y^2\to0 .
\]
Sıkıştırma teoremiyle limit $0$. (Kutupsal: $r^2(\cos^4\theta+\sin^4\theta)\le r^2$.)
\end{cevap}

\begin{soru}[\ekal Ek alıştırma 7]
\[
f(x,y)=\begin{cases}\dfrac{e^{x^2+y^2}-1}{x^2+y^2}, & (x,y)\neq(0,0)\\[6pt] k, & (x,y)=(0,0)\end{cases}
\]
fonksiyonu her yerde sürekli olacak şekilde $k$ sabitini bulunuz.
\end{soru}
\begin{cevap}
$(x,y)\neq(0,0)$'da $f$, sürekli fonksiyonların bölümü (payda $\neq0$) olduğundan süreklidir. Orijinde $u=x^2+y^2\to0^+$:
\[
\lim_{(x,y)\to(0,0)}f(x,y)=\lim_{u\to0^+}\frac{e^u-1}{u}=1
\]
(L'Hospital ya da $e^u$'nun $u=0$'daki türevinin $1$ olması). Süreklilik için $k=f(0,0)=1$.
\end{cevap}

\begin{soru}[\ekal Ek alıştırma 8]
$\displaystyle\lim_{(x,y,z)\to(0,0,0)}\frac{x^2+y^2-z^2}{x^2+y^2+z^2}$ limitini inceleyiniz.
\end{soru}
\begin{cevap}
$x$ ekseni ($y=z=0$): $\frac{x^2}{x^2}=1$. $z$ ekseni ($x=y=0$): $\frac{-z^2}{z^2}=-1$. İki farklı yol, iki farklı değer;
limit \textbf{yoktur}.
\end{cevap}

\begin{soru}[\ekal Ek alıştırma 9]
$g(x,y)=\dfrac{x^2y}{x^2+y^2}$ ($(x,y)\neq(0,0)$), $g(0,0)=0$ fonksiyonunun sürekli olduğu kümeyi bulunuz.
\end{soru}
\begin{cevap}
$(x,y)\neq(0,0)$'da rasyonel ve paydası sıfırdan farklı olduğundan süreklidir. Orijinde:
$\dfrac{x^2}{x^2+y^2}\le1$ olduğundan $|g(x,y)|\le|y|\to0=g(0,0)$. O hâlde $g$ \textbf{bütün düzlemde} süreklidir.
\end{cevap}
